\documentclass[11pt]{article}
\usepackage[a4paper,margin=25mm]{geometry}
\usepackage{amsmath,amssymb,amsthm,hyperref}
\newtheorem{lemma}{Lemma}\newtheorem{theorem}{Theorem}
\newcommand{\E}{\mathbb E}\newcommand{\Pol}{\operatorname{Pol}}
\title{Direct marked cancellation and range-sensitive endpoint errors}
\author{Complete auxiliary lemmas; independently reviewed}\date{30 September 2026}
\begin{document}\maketitle
Suppose nondecreasing comparison columns $p_1,\ldots,p_m$ and positive
probability row weights $w$ satisfy
$\E_w\prod_{j\in S}p_j=1/(|S|+1)$ for all subsets $S$.
Write $u_j=1-p_j$. Fix a marked column $i$, put $M=m-1$, choose
$\ell=\lfloor M/4\rfloor$, and put $N=M-\ell$.
For a uniformly chosen $\ell$-subset $B$ of the other $M$ columns,
let $R$ be its complement among those $M$ columns and set
\[
 A_B=\prod_{j\in B}p_j,\quad z_j=\ell u_j,\quad
 v_B=N^{-1}\sum_{j\in R}z_j,\quad
 x=M^{-1}\sum_{j\ne i}z_j,\quad h_i=z_i-x.
\]
All these quantities except $B,R$ are functions of the row.
For a univariate polynomial $P$, $\Pol_N(P)$ denotes its symmetric
multiaffine polarization, with $t^k$ replaced by
$e_k(z_1,\ldots,z_N)/\binom Nk$.

\begin{lemma}[Exact cancellation before scalar comparison]\label{identity}
For every polynomial $P$ of degree at most $N-1$,
\begin{equation}\label{exact}
 \E_w\E_B A_B\left[
 z_i\Pol_N(P)(z_R)-\Pol_N(tP(t))(z_R)\right]=0.
\end{equation}
For a fixed row and subset set
\[
 a_k=\frac{e_k((z_j-v_B)_{j\in R})}{\binom Nk}.
\]
Thus $a_0=1,a_1=0$. The same identity can be written
\begin{equation}\label{expanded}
 \E_w\E_B A_B(z_i-v_B)P(v_B)
 =\E_w\E_B A_B\sum_{k=2}^{\deg P+1}
 \left[\frac{(v_B-z_i)P^{(k)}(v_B)}{k!}
       +\frac{P^{(k-1)}(v_B)}{(k-1)!}\right]a_k,
\end{equation}
where derivatives above the degree are zero. The leading $k=2$ term
uses $a_2=-\operatorname{Var}_R(z)/(N-1)$.
\end{lemma}
\begin{proof}
For a monomial $P=t^r$, every term on the left of~\eqref{exact}
has $\ell$ distinct success factors and $r+1$ distinct failure
factors; the latter include the marked index on the first side and
are all in $R$ on the second. The mixed-moment identities, by expanding
the failure factors, give in both cases
\[
 \ell^{r+1}\int_0^1 t^\ell(1-t)^{r+1}dt.
\]
The index sets are disjoint and have total size at most $m$, since
$r+1\le N$. Linearity proves~\eqref{exact}.
The centered polarization formula is
\[
 \Pol_N(P)(z_R)=\sum_{k=0}^{\deg P}
                    \frac{P^{(k)}(v_B)}{k!}a_k.
\]
Also $(tP)^{(k)}(v_B)=v_BP^{(k)}(v_B)+kP^{(k-1)}(v_B)$.
Substitute these formulas into~\eqref{exact}, use $a_1=0$, and
rearrange. Finally $e_2(y)=-\sum_jy_j^2/2$ when $\sum_jy_j=0$.
\end{proof}

The point of~\eqref{expanded} is the cancellation of the common scalar
approximation error. There is no error term independent of the range
of the columns.

\begin{lemma}[A centered coefficient bound]\label{coefficient}
Let $H=\max_j z_j-\min_j z_j$, including the marked column if desired.
For every $2\le k\le N$,
\[
 |a_k|\le\left(C H\sqrt{k/N}\right)^k.
\]
\end{lemma}
\begin{proof}
Put $y_j=z_j-v_B$ and $\sigma^2=\sum_{j\in R}y_j^2\le NH^2$.
For every complex $t$, the elementary inequality
$|1+t y|\le\exp(\operatorname{Re}(ty)+|ty|^2/2)$ gives
\[
 \left|\prod_j(1+ty_j)\right|\le\exp(|t|^2\sigma^2/2),
\]
because $\sum_jy_j=0$. Cauchy's coefficient estimate at radius
$\sqrt{k}/\sigma$ yields
$|e_k(y)|\le(e\sigma^2/k)^{k/2}$.
Divide by $\binom Nk\ge(N/k)^k$. The case $\sigma=0$ is immediate.
\end{proof}

\begin{lemma}[Replacing the weighted subset mean]\label{subset}
Fix a row with $\min_{j\ne i}p_j\ge1/2$, and choose $B$ with
probability proportional to $A_B$. Suppose $H\le1$.
Then, with constants absolute for $M$ sufficiently large,
\[
 0\le\E_B(v_B-x)\le CH^2/M,\qquad
 \E_B(v_B-x)^2\le CH^2/M.
\]
Consequently, if $P$ is twice continuously differentiable on the
interval $I=[\min_jz_j,\max_jz_j]$, then
\begin{equation}\label{subseterror}
 \left|\E_B[(z_i-v_B)P(v_B)]-h_iP(x)\right|
 \le\frac{CH^2}{M}
       \left(\|P\|_{\infty,I}+\|P'\|_{\infty,I}
                                  +H\|P''\|_{\infty,I}\right).
\end{equation}
\end{lemma}
\begin{proof}
The exact success-weighted subset identity in
\path{../rational_designs/asymptotic/weighted_success_subset_moments.tex}
gives
\[
 \E_B(v_B-x)=\frac{\ell}{MN}
 \sum_{a<b}(p_a-p_b)^2
               \frac{e_{\ell-1}(p_{-[a,b]})}{e_\ell(p)}.
\]
The ratio is at most two, and
$\sum_{a<b}(p_a-p_b)^2
 =M\sum_a(p_a-\bar p)^2\le M^2H^2/\ell^2$.
Since $\ell,N$ are comparable to $M$, this proves the bias bound.
For the variance, subtract $\min_jz_j$ from every coefficient in the
selected sum; its variance is unchanged because $|B|=\ell$ is fixed.
All remaining coefficients are between $0$ and $H$. Pair inclusion
covariances are nonpositive, as proved in the cited note, so
\[
 \operatorname{Var}(v_B)\le N^{-2}\sum_{j\ne i}
                           (z_j-\min_a z_a)^2\le CH^2/M.
\]
Adding the squared bias proves the second-moment bound when $H\le1$.
For $g(v)=(z_i-v)P(v)$, on $I$ one has
\[
 |g'(v)|\le H|P'(v)|+|P(v)|,\qquad
 |g''(v)|\le H|P''(v)|+2|P'(v)|.
\]
Taylor's theorem at $x$, the bias, and the second moment give
\eqref{subseterror}.
\end{proof}

\begin{lemma}[Subgaussian subset shifts under every tilt]\label{subgaussian}
In the setting of Lemma~\ref{subset}, put $Y=v_B-x$. There is an
absolute constant $C$ such that
\[
 \log\E_B e^{t(Y-\E_BY)}\le CH^2t^2/M\qquad(t\in\mathbb R).
\]
For every fixed $C_0>0$, if $\tau\ge C_1H^2/M$ with a sufficiently
large $C_1=C_1(C_0)$, then
\[
 \E_B e^{C_0Y^2/\tau}\le C_{C_0},\qquad
 \E_B[Y^2e^{C_0Y^2/\tau}]\le C_{C_0}H^2/M.
\]
\end{lemma}
\begin{proof}
The exponential tilt of the weighted subset law by
$\exp(t\sum_{j\in B}z_j)$ is again a positive product-weight law on
fixed-size subsets, with weights $p_j e^{tz_j}$. Its pair inclusion
covariances are nonpositive by the same Newton-inequality proof; that
argument requires positivity, not an upper bound of one on the weights.
Subtracting $\min_jz_j$ from all coefficients leaves the variance
unchanged because the subset size is fixed. Thus, under every real tilt,
\[
 \operatorname{Var}_t\left(\sum_{j\in B}z_j\right)\le MH^2.
\]
The second derivative of its log moment-generating function is exactly
this tilted variance. Integrate the bound twice from zero and divide
by $N^2$, using $N\asymp M$, to obtain the displayed subgaussian bound
for $v_B$ centered at its own mean.
Lemma~\ref{subset} gives $|\E_BY|\le CH^2/M$. Since $H\le1$, its
squared mean is at most $CH^2/M$. The subgaussian tail estimate and
integration of the tail, or the elementary Gaussian integral identity
for $e^{cY^2}$, give the two final bounds when $C_1$ is large enough.
When $H=0$, $Y=0$ identically.
\end{proof}

\begin{lemma}[A localization-preserving subset comparison]\label{localized}
Retain those assumptions, take $0<\tau\le1$ with
$\tau\ge C_1H^2/M$, and fix a nonnegative envelope $A(x)$.
Suppose $P$ has, for $r=0,1,2$ and every $y\in I$,
\[
 |P^{(r)}(y)|\le C_0 A(x)\tau^{-r/2}
                       \exp(C_0(y-x)^2/\tau).
\]
Then
\begin{align*}
 |\E_BP(v_B)-P(x)|
   &\le \frac{C H^2}{M\tau}A(x),\\
 |\E_B[(z_i-v_B)P(v_B)]-h_iP(x)|
   &\le \frac{CH^2}{M}
           (1+\tau^{-1/2}+H\tau^{-1})A(x).
\end{align*}
The constant depends only on $C_0$.
\end{lemma}
\begin{proof}
Use Taylor's theorem at $x$. The first derivative term is bounded
by the bias $CH^2/M$ times the appropriate derivative at $x$.
For the integral second-order remainder, the displayed envelope at
$x+tY$ is bounded by its right side with $Y^2$ in the exponential.
Lemma~\ref{subgaussian} bounds its expectation, including the factor
$Y^2$, by $CH^2/M$. For the second assertion apply this argument to
$g(y)=(z_i-y)P(y)$ and use
\[
 |z_i-y|\le H,\quad
 g'=(z_i-y)P'-P,\quad g''=(z_i-y)P''-2P'.
\]
The stated bounds follow. This argument keeps the envelope at the
original mean $x$, instead of replacing it by a supremum on a
potentially much longer row-range interval.
\end{proof}

\begin{lemma}[Localized direct marked differential sum]\label{markedkernel}
Fix $0<a_0<L<\infty$, and suppose all row values $z_j$ and a kernel
center $a$ lie in $[a_0,L]$. Let $K_\tau$ be the Laguerre Poisson
kernel in \path{../localized_laguerre_kernel_derivatives.tex}.
For sufficiently small $0<\tau<\tau_0$ (depending only on the fixed interval),
if $H\le c$ and $H/\sqrt{N\tau}\le c$, then for every $d\le N$
\[
 \left|\sum_{k=2}^{d}
 \left[\frac{(v_B-z_i)\partial_v^kK_\tau(v_B,a)}{k!}
 +\frac{\partial_v^{k-1}K_\tau(v_B,a)}{(k-1)!}\right]a_k\right|
 \le C K_{4\tau}(v_B,a)
       \left(\frac{H^2}{N\sqrt\tau}+\frac{H^3}{N\tau}\right).
\]
The constants depend only on the fixed interval. Under the weighted
subset law, provided $\tau\ge C H^2/M$, the right side averages to
at most the same bound with $K_{C'\tau}(x,a)$ and $N$ replaced by $M$.
\end{lemma}
\begin{proof}
The derivative-order-$k$ part is the unmarked centered polarization
estimate from the cited note, multiplied by $|v_B-z_i|\le H$.
It is at most $C K_{4\tau}H^3/(N\tau)$ for $\tau\le1$.
For the derivative-order-$(k-1)$ part, combine that note's bound
\[
 \frac{|\partial_v^{k-1}K_\tau|}{(k-1)!}
 \le C K_{4\tau}
       \left[(C/\sqrt{(k-1)\tau})^{k-1}+C^{k-1}\right]
\]
with Lemma~\ref{coefficient}. The first series is bounded by
\[
 C K_{4\tau}\frac H{\sqrt N}
       \sum_{k=2}^{d}\sqrt{k}
                 (CH/\sqrt{N\tau})^{k-1}
 \le C K_{4\tau}\frac{H^2}{N\sqrt\tau}.
\]
For the second, use $(k/N)^{k/2}\le k/N$ and sum a geometric series
in $CH$, obtaining $C K_{4\tau}H^2/N$.
This proves the first assertion. The complex-shift bound in the
cited kernel note, applied to the real shift $v_B-x$, gives
\[
 K_{4\tau}(v_B,a)\le C K_{16\tau}(x,a)
                        \exp(C(v_B-x)^2/\tau)
\]
when the row range is sufficiently small. Lemma~\ref{subgaussian}
then proves the weighted averaging assertion.
\end{proof}

As with the corresponding unmarked kernel estimate, this lemma concerns
a finite differential sum. The exact polynomial identity
\eqref{expanded} can only use it after the spectral truncation of the
Poisson kernel has been quantitatively justified.

\paragraph{What this does and does not prove.}
Equations~\eqref{expanded} and~\eqref{subseterror} control the marked
functional $h_iP$ by errors quadratic in the row range $H$.
When all profiles agree, both errors vanish exactly. Combining this
with a positive localized polynomial test may improve the known
endpoint-profile rate. No such sharp rate is claimed here: one still
has to control localization, negative contributions outside a
one-sided deviation interval, and tails of the row measure.
\end{document}
